\section{Összefoglalás}

A jelen dolgozat egy valós idejű, mesterséges intelligencián alapuló 3D nyomtatási hibadetektor rendszer
tervezését és megvalósítását mutatta be, amelynek célja az FDM technológiával üzemelő nyomtatófarmok
automatizált felügyelete.

A munka első szakaszában az iparági igények és a kapcsolódó szakirodalom elemzésével meghatároztuk azokat
a leggyakoribb és gazdaságilag legjelentősebb hibatípusokat -- spagetti, nem tapadó réteg, réteg-eltolódás,
vetemedés, szálazás, alul- és felülextrúdálás, fúvóka-dugulás, idegen tárgy --, amelyek automatikus
detektálása a legtöbb értéket hozza a nyomtatógazdálkodásban.

Az adatgyűjtési fázisban három nyilvánosan elérhető adathalmaz, automatizált webes képletöltés és 397 saját
készítésű felvétel kombinációjából \textbf{2422 annotált képből} álló adathalmazt állítottunk össze.
Az annotálás CVAT eszközzel, kilenc osztályra, bounding-box formátumban történt. Az osztályegyensúly-hiányt
oversampling technikával, a vizuális diverzitást pedig egy háromrétegű augmentációs pipeline-nal --
alaptranszformációk, Smart MixUp és mozaik -- korrigáltuk, amelynek eredményeként a tanítókészlet közel
négyszeresére bővült.

A modellezési szakaszban a YOLOv8x architektúrát alkalmaztuk, kétfázisú tanítási stratégiával.
Az első fázis 960 képpontos bemeneti felbontáson általánosabb jellemzők megtanulását célozta, míg a
második fázis 1280 képpontos felbontáson, az első fázis legjobb súlyaiból kiindulva, finomhangolással
tovább javított a részletgazdag hibaminták felismerésén. A betanított modell végső teljesítménymutatói:

\begin{table}[H]
\centering
\caption{A betanított YOLOv8x modell összesített teljesítménye}
\renewcommand{\arraystretch}{1.4}
\begin{tabular}{|l|c|}
    \hline
    \textbf{Metrika} & \textbf{Érték} \\
    \hline
    Pontosság (Precision) & $0{,}937$ \\
    \hline
    Visszahívás (Recall)  & $0{,}836$ \\
    \hline
    mAP@0.5               & $0{,}892$ \\
    \hline
    mAP@0.5:0.95          & $0{,}677$ \\
    \hline
    Csúcs mAP@0.5         & $0{,}901$ \\
    \hline
    Csúcs F1              & $0{,}88$ \\
    \hline
\end{tabular}
\end{table}

A rendszerintegrációs szakaszban egy edge--cloud hibrid architektúra valósult meg. A helyszíni komponenst
egy Raspberry Pi 4B adja, amelyen három webkamera képkockáit egy \texttt{frame-extractor} szolgáltatás
gyűjti és publikálja Google Cloud Pub/Sub csatornára. A felhőoldali pipeline egy Dispatcher Cloud Function,
egy Vertex AI-n futó \texttt{judge} konténer (YOLOv8 inferenciával), egy Alert-Manager Cloud Function,
valamint egy Firestore-alapú valós idejű webes dashboard elemeiből épül fel. Az infrastruktúra
Terraformmal kezelt, a modellfrissítések CI/CD workflow-n keresztül történnek -- emberi beavatkozás nélkül.

A rendszer három kameraállásból, másodpercenként hat képkockát dolgoz fel, és magas prioritású hiba
észlelésekor e-mail értesítést küld a kezelőnek, 60 másodperces kameránkénti cooldown-nal a riasztásözön
elkerülése érdekében.

% -------------------------------------------------------------------
\section{Jövőbeli irányok}

A jelenlegi megvalósítás több ponton is tovább fejleszthető:

\begin{itemize}
    \item \textbf{Konfidenciaküszöb finomhangolása.} A jelenlegi 35\%-os küszöb empirikusan lett
    meghatározva. Éles nyomtatási körülmények között szisztematikus mérésekkel -- precision--recall
    kompromisszum elemzéssel -- pontosabb, üzemspecifikus érték határozható meg.

    \item \textbf{Állandó Cloudflare alagút.} Az ideiglenes \textit{quick tunnel} minden Pi-újraindításkor
    új URL-t kap, ami megszakítja a dashboard stream-kapcsolatát. Éles üzemeltetéshez \textit{named tunnel}
    konfigurációra van szükség.

    \item \textbf{Vertex AI scale-to-zero.} A végpont jelenlegi minimálpéldánya napi $\approx$37 USD
    üzemelési költséggel jár. Alacsony intenzitású monitorozási forgatókönyvekben a \textit{scale-to-zero}
    beállítás jelentős megtakarítást hozhat a hideginduló késleltetés elfogadható növelése árán.

    \item \textbf{Bővített osztálykészlet.} Az adathalmaz elsősorban az FDM technológia legelterjedtebb
    hibatípusait fedi le. Speciálisabb hibák -- például fúvóka-kopás, rétegleválás, hőmérsékleti
    anomáliák -- bevonásával a rendszer alkalmazhatósági köre bővíthető.

    \item \textbf{Modell könnyűsúlyúsítása.} A YOLOv8x nagy kapacitású modell; edge eszközön történő
    futtatáshoz YOLOv8n vagy YOLOv8s változatra való átállás, illetve kvantálás csökkenthetné
    a felhőinfrastruktúra igénybevételét és a késleltetést.

    \item \textbf{Aktív tanulás.} A rendszer üzemelés közben gyűjti a valós nyomtatási képeket.
    Egy aktív tanulási ciklus bevezetésével az operátor által megerősített vagy korrigált detekciók
    visszacsatolhatók az adathalmazba, folyamatosan javítva a modell teljesítményét.

    \item \textbf{Nyomtatóvezérlő integráció.} Jelenleg a rendszer csak értesít; a következő lépés
    a detektált hiba súlyosságától függő automatikus beavatkozás -- például a nyomtatás szüneteltetése
    vagy leállítása -- megvalósítása lenne OctoPrint vagy Klipper API-n keresztül.
\end{itemize}
